\expandafter\gdef\csname SNside1\endcsname{\node[anchor=north west,inner sep=0pt,text width=48mm,font=\SideFont,align=left] (SNsidebox) at ([xshift=138mm,yshift=-375.69pt]current page.north west) {\textcolor{SNink}{\textbf{最小反例 · 数量相同仍不合法}}\par\vspace{1.5mm}顺序输入a、b、c，要先输出c就必须把三者全部入栈；弹出c后栈顶为b，接着要求a便失败。输出c、a、b与输入含相同元素，却违反栈顶约束。\par\smallskip\begin{tikzpicture}[x=1mm,y=1mm,>=Stealth,node distance=5mm]
\node[draw=SNink!50,fill=SNblue,rounded corners=1mm,text width=34mm,align=center,inner sep=1.2mm,font=\fontsize{7.5}{9}\selectfont] (m0) at (0,0) {底a};
\node[draw=SNink!50,fill=SNblue,rounded corners=1mm,text width=34mm,align=center,inner sep=1.2mm,font=\fontsize{7.5}{9}\selectfont] (m1) at (0,-11) {顶b};
\node[draw=SNink!50,fill=SNblue,rounded corners=1mm,text width=34mm,align=center,inner sep=1.2mm,font=\fontsize{7.5}{9}\selectfont] (m2) at (0,-22) {要求a};
\draw[->,SNink] (m0) -- (m1);
\draw[->,SNink] (m1) -- (m2);
\end{tikzpicture}
\par{\fontsize{7.5}{9}\selectfont c已弹出，b挡在a上方}};\pgfextracty{\SNsideheight}{\pgfpointdiff{\pgfpointanchor{SNsidebox}{south}}{\pgfpointanchor{SNsidebox}{north}}}\typeout{SNHEIGHT:1:0:\the\SNsideheight}}
\expandafter\gdef\csname SNside2\endcsname{\node[anchor=north west,inner sep=0pt,text width=48mm,font=\SideFont,align=left] (SNsidebox) at ([xshift=138mm,yshift=-95.00pt]current page.north west) {\textcolor{SNink}{\textbf{错项排除 · 共享栈只在相邻时满}}\par\vspace{1.5mm}两栈从数组两端向中间增长，未用单元恰位于两栈顶之间。两栈顶相邻已经没有空位；等到相等才判满，就已让两个栈占用同一个位置。};\pgfextracty{\SNsideheight}{\pgfpointdiff{\pgfpointanchor{SNsidebox}{south}}{\pgfpointanchor{SNsidebox}{north}}}\typeout{SNHEIGHT:2:2:\the\SNsideheight}}
\expandafter\gdef\csname SNside3\endcsname{\node[anchor=north west,inner sep=0pt,text width=48mm,font=\SideFont,align=left] (SNsidebox) at ([xshift=138mm,yshift=-95.00pt]current page.north west) {\textcolor{SNink}{\textbf{辨析 · 循环队列先定约定}}\par\vspace{1.5mm}牺牲一个单元时，空为front=rear，满为(rear+1)\%MaxSize=front；其他设计的条件需另算。};\pgfextracty{\SNsideheight}{\pgfpointdiff{\pgfpointanchor{SNsidebox}{south}}{\pgfpointanchor{SNsidebox}{north}}}\typeout{SNHEIGHT:3:21:\the\SNsideheight}}
\expandafter\gdef\csname SNside4\endcsname{\node[anchor=north west,inner sep=0pt,text width=48mm,font=\SideFont,align=left] (SNsidebox) at ([xshift=138mm,yshift=-95.00pt]current page.north west) {\textcolor{SNink}{\textbf{不变量 · 尾指针不能悬空}}\par\vspace{1.5mm}带头结点链队列始终让rear指向最后一个数据结点，空时指向头结点。删去唯一数据结点后，rear必须回到front；否则下次入队会经已释放地址写内存。};\pgfextracty{\SNsideheight}{\pgfpointdiff{\pgfpointanchor{SNsidebox}{south}}{\pgfpointanchor{SNsidebox}{north}}}\typeout{SNHEIGHT:4:4:\the\SNsideheight}}
\expandafter\gdef\csname SNside5\endcsname{\node[anchor=north west,inner sep=0pt,text width=48mm,font=\SideFont,align=left] (SNsidebox) at ([xshift=138mm,yshift=-95.00pt]current page.north west) {\textcolor{SNink}{\textbf{错项排除 · 三对角首尾行少一个}}\par\vspace{1.5mm}n至少为2时，首行和末行各只保存2个位置，中间行才各3个。若把第i行之前全部按3个计，会从第一行就多算；先列2+3(i\ensuremath{-}2)再加本行偏移。};\pgfextracty{\SNsideheight}{\pgfpointdiff{\pgfpointanchor{SNsidebox}{south}}{\pgfpointanchor{SNsidebox}{north}}}\typeout{SNHEIGHT:5:8:\the\SNsideheight}}
\expandafter\gdef\csname SNside6\endcsname{\node[anchor=north west,inner sep=0pt,text width=48mm,font=\SideFont,align=left] (SNsidebox) at ([xshift=138mm,yshift=-95.00pt]current page.north west) {\textcolor{SNink}{\textbf{最小反例 · 括号计数看不到嵌套}}\par\vspace{1.5mm}({[}){]}中每一类左右括号数都相等，但第三个字符要求圆括号配对，栈顶却是方括号。配对必须遵守最近未完成对象先结束，不能在扫描后任意重排。};\pgfextracty{\SNsideheight}{\pgfpointdiff{\pgfpointanchor{SNsidebox}{south}}{\pgfpointanchor{SNsidebox}{north}}}\typeout{SNHEIGHT:6:9:\the\SNsideheight}}
\expandafter\gdef\csname SNside7\endcsname{\node[anchor=north west,inner sep=0pt,text width=48mm,font=\SideFont,align=left] (SNsidebox) at ([xshift=138mm,yshift=-95.00pt]current page.north west) {\textcolor{SNink}{\textbf{辨析 · 调用次数与同时深度}}\par\vspace{1.5mm}Fibonacci调用树的结点总数反映总调用次数；栈空间取同时尚未返回的调用链长度。};\pgfextracty{\SNsideheight}{\pgfpointdiff{\pgfpointanchor{SNsidebox}{south}}{\pgfpointanchor{SNsidebox}{north}}}\typeout{SNHEIGHT:7:22:\the\SNsideheight}}
\expandafter\gdef\csname SNside8\endcsname{\node[anchor=north west,inner sep=0pt,text width=48mm,font=\SideFont,align=left] (SNsidebox) at ([xshift=138mm,yshift=-95.00pt]current page.north west) {\textcolor{SNink}{\textbf{不变量 · top含义决定增减位置}}\par\vspace{1.5mm}top指栈顶时，空状态为\ensuremath{-}1，入栈先加再写；top指下一空位时，空状态为0，入栈先写再加。用空栈和单元素两幅状态检查，比只背push语句更能应对改约定题。};\pgfextracty{\SNsideheight}{\pgfpointdiff{\pgfpointanchor{SNsidebox}{south}}{\pgfpointanchor{SNsidebox}{north}}}\typeout{SNHEIGHT:8:1:\the\SNsideheight}}
\expandafter\gdef\csname SNside9\endcsname{\node[anchor=north west,inner sep=0pt,text width=48mm,font=\SideFont,align=left] (SNsidebox) at ([xshift=138mm,yshift=-95.00pt]current page.north west) {\textcolor{SNink}{\textbf{不变量 · 先换到已保存半边}}\par\vspace{1.5mm}对称矩阵访问未存的半边时，先交换i、j，再按保存方向计数。坐标转换解决的是表示对应；数组从0还是1起始解决的是编号，两步不能互相替代。};\pgfextracty{\SNsideheight}{\pgfpointdiff{\pgfpointanchor{SNsidebox}{south}}{\pgfpointanchor{SNsidebox}{north}}}\typeout{SNHEIGHT:9:7:\the\SNsideheight}}
\expandafter\gdef\csname SNside10\endcsname{\node[anchor=north west,inner sep=0pt,text width=48mm,font=\SideFont,align=left] (SNsidebox) at ([xshift=138mm,yshift=-95.00pt]current page.north west) {\textcolor{SNink}{\textbf{错项排除 · FIFO只是具体策略}}\par\vspace{1.5mm}队列ADT规定先进先出，操作系统的``就绪队列''名称却不保证所有调度算法都按FIFO选择。先来先服务可借FIFO说明，优先级等策略需要额外排序或选择规则。};\pgfextracty{\SNsideheight}{\pgfpointdiff{\pgfpointanchor{SNsidebox}{south}}{\pgfpointanchor{SNsidebox}{north}}}\typeout{SNHEIGHT:10:12:\the\SNsideheight}}
\expandafter\gdef\csname SNside11\endcsname{\node[anchor=north west,inner sep=0pt,text width=48mm,font=\SideFont,align=left] (SNsidebox) at ([xshift=138mm,yshift=-95.00pt]current page.north west) {\textcolor{SNink}{\textbf{辨析 · 栈顶选在链头}}\par\vspace{1.5mm}单链表头插头删只改常数条链接。若从单链表尾部弹出，还要先寻找前驱，代价会增加。};\pgfextracty{\SNsideheight}{\pgfpointdiff{\pgfpointanchor{SNsidebox}{south}}{\pgfpointanchor{SNsidebox}{north}}}\typeout{SNHEIGHT:11:13:\the\SNsideheight}}
\expandafter\gdef\csname SNside12\endcsname{\node[anchor=north west,inner sep=0pt,text width=48mm,font=\SideFont,align=left] (SNsidebox) at ([xshift=138mm,yshift=-95.00pt]current page.north west) {\textcolor{SNink}{\textbf{辨析 · 最后一删要回尾}}\par\vspace{1.5mm}删除唯一数据结点后，rear必须退回front；二者共同指向仍然有效的头结点。};\pgfextracty{\SNsideheight}{\pgfpointdiff{\pgfpointanchor{SNsidebox}{south}}{\pgfpointanchor{SNsidebox}{north}}}\typeout{SNHEIGHT:12:14:\the\SNsideheight}}
\expandafter\gdef\csname SNside13\endcsname{\node[anchor=north west,inner sep=0pt,text width=48mm,font=\SideFont,align=left] (SNsidebox) at ([xshift=138mm,yshift=-95.00pt]current page.north west) {\textcolor{SNink}{\textbf{辨析 · 共享不改变次序}}\par\vspace{1.5mm}共享栈只共享空闲容量，两栈仍各自后进先出。top1-top0=1表示中间已无空格。};\pgfextracty{\SNsideheight}{\pgfpointdiff{\pgfpointanchor{SNsidebox}{south}}{\pgfpointanchor{SNsidebox}{north}}}\typeout{SNHEIGHT:13:15:\the\SNsideheight}}
\expandafter\gdef\csname SNside14\endcsname{\node[anchor=north west,inner sep=0pt,text width=48mm,font=\SideFont,align=left] (SNsidebox) at ([xshift=138mm,yshift=-95.00pt]current page.north west) {\textcolor{SNink}{\textbf{辨析 · 相同优先级看结合性}}\par\vspace{1.5mm}乘除和加减通常左结合，中缀转后缀遇同级运算符时先弹栈顶。新运算符不能抢先。};\pgfextracty{\SNsideheight}{\pgfpointdiff{\pgfpointanchor{SNsidebox}{south}}{\pgfpointanchor{SNsidebox}{north}}}\typeout{SNHEIGHT:14:17:\the\SNsideheight}}
\expandafter\gdef\csname SNside15\endcsname{\node[anchor=north west,inner sep=0pt,text width=48mm,font=\SideFont,align=left] (SNsidebox) at ([xshift=138mm,yshift=-95.00pt]current page.north west) {\textcolor{SNink}{\textbf{辨析 · 三角区与常量区}}\par\vspace{1.5mm}对称矩阵只存半边；三角矩阵还需存另一半的公共常量。两者存储单元数相差1。};\pgfextracty{\SNsideheight}{\pgfpointdiff{\pgfpointanchor{SNsidebox}{south}}{\pgfpointanchor{SNsidebox}{north}}}\typeout{SNHEIGHT:15:18:\the\SNsideheight}}
\expandafter\gdef\csname SNside16\endcsname{\node[anchor=north west,inner sep=0pt,text width=48mm,font=\SideFont,align=left] (SNsidebox) at ([xshift=138mm,yshift=-95.00pt]current page.north west) {\textcolor{SNink}{\textbf{自查提示 · 队列的满空条件}}\par\vspace{1.5mm}写清采用牺牲一个单元、计数器还是标志位，才能计算容量并判断队满。};\pgfextracty{\SNsideheight}{\pgfpointdiff{\pgfpointanchor{SNsidebox}{south}}{\pgfpointanchor{SNsidebox}{north}}}\typeout{SNHEIGHT:16:20:\the\SNsideheight}}
\expandafter\gdef\csname SNside17\endcsname{\node[anchor=north west,inner sep=0pt,text width=48mm,font=\SideFont,align=left] (SNsidebox) at ([xshift=138mm,yshift=-95.00pt]current page.north west) {\textcolor{SNink}{\textbf{自查提示 · 队列的满空条件}}\par\vspace{1.5mm}写清采用牺牲一个单元、计数器还是标志位，才能计算容量并判断队满。};\pgfextracty{\SNsideheight}{\pgfpointdiff{\pgfpointanchor{SNsidebox}{south}}{\pgfpointanchor{SNsidebox}{north}}}\typeout{SNHEIGHT:17:20:\the\SNsideheight}}
\AddToShipoutPictureFG{\begin{tikzpicture}[remember picture,overlay]\ifcsname SNside\thepage\endcsname\csname SNside\thepage\endcsname\fi\end{tikzpicture}}